\documentclass[10pt,a4paper]{article}
\usepackage[utf8]{inputenc}
\usepackage[T1]{fontenc}
\usepackage{lmodern}
\usepackage{amsmath,amssymb}
\usepackage{geometry}
\usepackage{array,booktabs,tabularx}
\usepackage[most]{tcolorbox}
\usepackage{enumitem}
\usepackage{hyperref}

\geometry{margin=1.8cm, top=1.6cm, bottom=1.6cm}
\setlength{\parindent}{0pt}
\setlength{\parskip}{0.45em}
\pagestyle{empty}

\newtcolorbox{exambox}[1][]{
    colback=white,
    colframe=black!75,
    boxrule=0.8pt,
    arc=2pt,
    left=9pt,
    right=9pt,
    top=6pt,
    bottom=6pt,
    fonttitle=\bfseries\small,
    coltitle=black,
    colbacktitle=black!8,
    #1
}

\newtcolorbox{heuristicbox}[1][]{
    colback=blue!3!white,
    colframe=blue!55!black,
    boxrule=0.8pt,
    arc=2pt,
    left=9pt,
    right=9pt,
    top=6pt,
    bottom=6pt,
    fonttitle=\bfseries\small,
    #1
}

\begin{document}

\begin{center}
    {\Large \textbf{Questão Original da Prova de Matemática Discreta}}\\[0.2em]
    {\normalsize Universidade Federal de Minas Gerais (DCC / ICEx / UFMG)}\\[0.15em]
    {\footnotesize Baseada em Kenneth H. Rosen, \textit{Discrete Mathematics and Its Applications}, Cap. 7, Seção 7.1, Exercícios 36 e 37}
\end{center}

\vspace{0.2em}

\textbf{4. (20 pontos)}

\textbf{(a) O que é mais provável: sair um total de 8 quando dois dados são lançados ou sair um total de 8 quando três dados são lançados?}

\begin{exambox}[title={Resolução combinatória por contagem direta (conforme folha de prova)}]
\renewcommand{\arraystretch}{1.2}
\begin{tabularx}{\textwidth}{@{}p{0.37\textwidth}|X@{}}
\textbf{Lançamento de 2 dados} & \textbf{Lançamento de 3 dados} \\
\hline
\textbf{Soma 8} ($5$ pares ordenados): &
\textbf{Soma 8} ($21$ trincas ordenadas): \\
$(2,6),\; (6,2),$ &
$(1,1,6),\; (1,6,1),\; (6,1,1),\; (1,2,5),\; (1,5,2),\; (2,1,5),\; (2,5,1),$ \\
$(3,5),\; (5,3),$ &
$(5,1,2),\; (5,2,1),\; (1,3,4),\; (1,4,3),\; (3,1,4),\; (3,4,1),\; (4,1,3),$ \\
$(4,4)$ &
$(4,3,1),\; (2,2,4),\; (2,4,2),\; (4,2,2),\; (2,3,3),\; (3,2,3),\; (3,3,2)$ \\[0.35em]
Probabilidade: $\displaystyle P(S_2 = 8) = \frac{5}{36} = \frac{30}{216}$ &
Probabilidade: $\displaystyle P(S_3 = 8) = \frac{21}{216}$
\end{tabularx}

\vspace{0.35em}
\hrule
\vspace{0.35em}
\centering
Comparação: \quad
$\displaystyle \boxed{\frac{30}{216}} > \frac{21}{216} \quad \Longrightarrow \quad$
\textbf{Mais provável quando 2 dados são lançados.}
\end{exambox}

\vspace{0.25em}

\textbf{(b) O que é mais provável: sair um total de 9 quando dois dados são lançados ou sair um total de 9 quando três dados são lançados?}

\begin{exambox}[title={Resolução combinatória por contagem direta (conforme folha de prova)}]
\textbf{2 dados} ($4$ pares ordenados): \quad
$(3,6),\; (4,5),\; (5,4),\; (6,3) \quad \longrightarrow \quad \displaystyle P(S_2 = 9) = \frac{4}{36} = \frac{24}{216}$

\vspace{0.3em}
\textbf{3 dados} ($25$ trincas ordenadas, organizadas pelo resultado do primeiro dado):
\vspace{-0.4em}
\begin{align*}
d_1 = 1: &\quad (1,2,6),\; (1,3,5),\; (1,4,4),\; (1,5,3),\; (1,6,2) && \text{($5$ casos)}\\
d_1 = 2: &\quad (2,1,6),\; (2,2,5),\; (2,3,4),\; (2,4,3),\; (2,5,2),\; (2,6,1) && \text{($6$ casos)}\\
d_1 = 3: &\quad (3,1,5),\; (3,2,4),\; (3,3,3),\; (3,4,2),\; (3,5,1) && \text{($5$ casos)}\\
d_1 = 4: &\quad (4,1,4),\; (4,2,3),\; (4,3,2),\; (4,4,1) && \text{($4$ casos)}\\
d_1 = 5: &\quad (5,1,3),\; (5,2,2),\; (5,3,1) && \text{($3$ casos)}\\
d_1 = 6: &\quad (6,1,2),\; (6,2,1) && \text{($2$ casos)}
\end{align*}
\vspace{-0.7em}
\hrule
\vspace{0.35em}
\centering
$\displaystyle P(S_3 = 9) = \frac{5 + 6 + 5 + 4 + 3 + 2}{6^3} = \frac{25}{216}
\qquad \Longrightarrow \qquad
\frac{25}{216} > \frac{4}{36} = \frac{24}{216} \quad \Longrightarrow \quad \boxed{\textbf{3 dados}}$
\end{exambox}

\vspace{0.25em}

\begin{heuristicbox}[title={Resolução alternativa apresentada em prova (Heurística da Distância até a Média)}]
Em vez de enumerar todas as tuplas, um estudante comparou a distância absoluta de cada soma alvo até a esperança (média) de cada lançamento, sabendo que a média de um dado honesto de $6$ faces é $\mu_1 = 3{,}5$:
\[
E[S_2] = 2 \times 3{,}5 = 7{,}0 \qquad \text{e} \qquad E[S_3] = 3 \times 3{,}5 = 10{,}5
\]
\vspace{-0.6em}
\begin{itemize}[leftmargin=1.5em, itemsep=0.2em, topsep=0.1em]
    \item \textbf{Item (a), soma $k = 8$:} \quad
    $|8 - E[S_2]| = |8 - 7{,}0| = 1{,}0 < |8 - E[S_3]| = |8 - 10{,}5| = 2{,}5 \;\Longrightarrow\;$ \textbf{2 dados.}
    \item \textbf{Item (b), soma $k = 9$:} \quad
    $|9 - E[S_3]| = |9 - 10{,}5| = 1{,}5 < |9 - E[S_2]| = |9 - 7{,}0| = 2{,}0 \;\Longrightarrow\;$ \textbf{3 dados.}
\end{itemize}
\vspace{0.1em}
\textbf{Por que funcionou na prova, mas falha no caso geral?} O ponto médio aritmético entre $7{,}0$ e $10{,}5$ é $8{,}75$, que cai entre os inteiros $8$ e $9$, exatamente no mesmo intervalo que o ponto de virada contínuo real ($x^* \approx 8{,}98$). Por essa coincidência numérica, a heurística acerta os dois itens da avaliação, embora falhe quando o número de dados se afasta (por exemplo, entre $2$ e $5$ dados na soma $12$, ou entre $2$ e $6$ dados na soma $13$).
\end{heuristicbox}

\end{document}
